\scanpage{48}\mspage{54}
On dit alors que $f$ est \uline{\deleted{\unsure{co}}homotopiquement propre}. La \deleted{\unsure{condition}} \inserted{notion} duale s'exprime, soit par la condition que les foncteurs induits par $f$
\[
 X_{y_0}\longrightarrow {}_{y_0\backslash}Y
\]
soient des hotéquivalences strictes i.e. qu'ils soient cohomotopiquement cofinaux, soit que les fibres des morphismes induits sur les ``catégories localisées'' (catégories des générisations)
\[
 X_{/x_0}\longrightarrow Y_{/f(x_0)}
\]
sont contractiles. (N.B., On utilise ici qu'une catégorie $Z$ est contractile ssi $Z^\circ$ l'est.) On dit alors que $f$ est \uline{homotopiquement lisse}

\uline{Prop} a) La notion de foncteur propre \inserted{(\unsure{resp.|ou} lisse)} est stable par changement de base, par composition. \inserted{Les isomorphismes.} \deleted{Plus gén. les \unsure{si} $gf$ est propre, et $g$ propre, alors $f$ est propre ; i.e. tout $S$-foncteur entre catégories $X,Y$ sur $S$ propres sur $S$, est propre.}

% slanted
\margin{Attention, une équiv. de cat. n'est pas propre !}
\deleted{b) L'énoncé dual pour foncteurs lisses. Notamment, si $gf$ est lisse et \unsure{l'équivalence de catégories}}
\inserted{\unsure{Mais oui|Mais} les immers. fermées, sont propres. L'image d'un morphisme propre est une \deleted{\unsure{partie de}} sous-catégorie fermée (\unsure{non pleine})$^*$}
% slanted
\margin{Si $X'\hookrightarrow X$ inclusion d'une sous-catégorie de $X$, elle \deleted{\unsure{est}} est propre \unsure{ssi|$\Leftrightarrow$} \deleted{\unsure{si}} c'est une \unsure{immers.|immersion} fermée}

\deleted{\unsure{Cor.}} \uline{NB} le foncteur $X\to e$ est tjrs propre et lisse, car les cat ${}_{x_0\backslash}X$ et $X_{/x_0}$ sont contractiles (ayant objet initial ou objet final). Donc tt foncteur de projection $X\times Y\xrightarrow[\mathrm{pr}_{\unsure{1|,}}]{}Y$ est à la fois propre et lisse.
% archive sheet number: 47
